\chapter{חלופות ובחירת פתרון}\label{ch:alt}

נבחנו ארבע חלופות ממשיות. כל אחת מייצגת פילוסופיה אחרת של פתרון, וכל אחת נבחרה בפועל על ידי שחקן כלשהו בשוק.

\section{תיאור החלופות}
\begin{description}[leftmargin=0pt,labelsep=6pt]
\item[חלופה א׳: גלאי בלבד.] חיבור מערכת ניהול התוכן לממשק \en{API} מסחרי לזיהוי בינה מלאכותית והצגת הציון לעורך. זה הפתרון הזול והמהיר ביותר ליישום, אבל הוא נשען על הרכיב הפחות אמין (פרק~\ref{ch:lit}) ואינו מטפל ברובד האתי כלל.
\item[חלופה ב׳: מקור בלבד.] מדיניות ״רק תמונות חתומות״: המערכת מקבלת רק תמונות עם תעודת \en{C2PA} תקפה. הפתרון ודאי כשהוא חל, אבל פוסל כ-\num{99\%} מהתמונות הקיימות~\cite{c2pa2026}, כולל תיעוד אזרחי חשוב מזירות אסון ומלחמה.
\item[חלופה ג׳: בדיקה אנושית במיקור חוץ.] העברת כל תמונה חשודה לגוף בדיקת עובדות חיצוני. הבדיקה איכותית, אבל איטית (שעות), יקרה, ואינה עומדת בקנה מידה של מיליוני תמונות.
\item[חלופה ד׳: היברידית (\en{VeriLens}).] שלוש שכבות אימות אוטומטיות (מקור, זיהוי, הקשר), ציון מכויל עם מצב ״לא ודאי״, החלטה אנושית מתועדת, תיוג לקורא, ותקן אתי שמגדיר את הנוהל.
\end{description}

\section{קריטריונים ומשקלות}
המשקלות נקבעו בסבב של ״ועדת המומחים״ (שבעה תפקידי מומחה שהפעלנו כסוכני בינה מלאכותית, ראו סעיף~\ref{sec:committee}) לפני שהחלופות נוקדו, כדי למנוע התאמה בדיעבד (טבלה~\ref{tab:decision}).

\begin{table}[!htbp]\centering
\caption{מטריצת החלטה משוקללת (ציון \num{1} עד \num{5} לכל קריטריון)}\label{tab:decision}
\begin{tabularx}{\linewidth}{@{}Y c c c c c@{}}
\toprule
\textbf{קריטריון} & \textbf{משקל} & \textbf{א׳ גלאי} & \textbf{ב׳ מקור} & \textbf{ג׳ אנושי} & \textbf{ד׳ היברידי} \\
\midrule
דיוק ההחלטה & \num{25\%} & \num{2} & \num{5} & \num{4} & \num{4} \\
כיסוי (אחוז התמונות שאפשר לבדוק) & \num{20\%} & \num{5} & \num{1} & \num{3} & \num{5} \\
מהירות & \num{15\%} & \num{5} & \num{4} & \num{1} & \num{4} \\
אחריותיות ותיעוד (רובד אתי) & \num{15\%} & \num{1} & \num{3} & \num{4} & \num{5} \\
יכולת גדילה & \num{10\%} & \num{5} & \num{4} & \num{1} & \num{4} \\
עלות ליחידה (\num{5} = זול) & \num{10\%} & \num{4} & \num{4} & \num{1} & \num{2} \\
התאמה לרגולציה & \num{5\%} & \num{2} & \num{4} & \num{3} & \num{5} \\
\midrule
\textbf{ציון משוקלל (מתוך \num{100})} & & \textbf{68} & \textbf{70} & \textbf{54} & \textbf{84} \\
\bottomrule
\end{tabularx}
\end{table}

ראו איור~\ref{fig:altscore}.


\begin{figure}[!htbp]
\begin{ltrdiagram}
% altscore: TikZ / pgfplots source, built inside ltrdiagram (LTR)
\begin{tikzpicture}
\begin{axis}[
  xbar, width=11cm, height=4.6cm, xmin=0, xmax=100,
  symbolic y coords={C,A,B,D}, ytick={C,A,B,D},
  yticklabels={\he{ג׳ בדיקה אנושית},\he{א׳ גלאי בלבד},\he{ב׳ מקור בלבד},\he{ד׳ היברידי, נבחר}},
  yticklabel style={font=\sffamily\small},
  xlabel={\he{ציון משוקלל, מתוך 100}}, xlabel style={font=\sffamily\small},
  nodes near coords, nodes near coords style={font=\sffamily\small},
  bar width=11pt, axis x line*=bottom, axis y line*=left, xmajorgrids, grid style={gray!25},
  enlarge y limits=0.2]
\addplot[fill=acc!35,draw=acc,bar shift=0pt] coordinates {(54,C) (68,A) (70,B)};
\addplot[fill=acc2,draw=acc2,bar shift=0pt] coordinates {(84,D)};
\end{axis}
\end{tikzpicture}

\end{ltrdiagram}
\caption{ציון משוקלל לכל חלופה: ההיברידית מובילה ב-\num{14} נקודות}\label{fig:altscore}
\end{figure}

\section{ניתוח רגישות}
כדי לוודא שהבחירה אינה תלויה במשקלות שרירותיים, נבדקו שלושה תרחישים קיצוניים:
\begin{itemize}
\item \textbf{״העלות קובעת״} (משקל העלות הוכפל ל-\num{20\%} על חשבון הדיוק): ד׳ = \num{80}, א׳ = \num{72}. ההיברידית עדיין מובילה.
\item \textbf{״הדיוק קובע״} (דיוק \num{40\%}, כיסוי \num{10\%}, מהירות \num{10\%}): ד׳ = \num{82}, ב׳ = \num{79}. הפער מצטמצם, אבל הסדר נשמר.
\item \textbf{״ללא רובד אתי״} (משקל האחריותיות אופס והשאר נורמלו): ד׳ = \num{81}, א׳ = \num{77}. הפער מצטמצם מ-\num{16} ל-\num{4} נקודות, ומכאן לקח חשוב: בלי הרובד האתי, היתרון של \en{VeriLens} על גלאי זול כמעט נעלם. התקן הוא חלק מהותי מהמוצר, לא תוספת.
\end{itemize}

\section{הבחירה ונימוקה}
\begin{keybox}[הוחלט]
חלופה ד׳: היברידית. היא היחידה שמשיגה כיסוי מלא בלי לוותר על ודאות כשהיא קיימת, שומרת את ההחלטה בידי אדם כפי שדורשת אמנת פריז~\cite{paris2023}, ומייצרת את התיעוד שהרגולציה דורשת. המחיר הוא עלות פיתוח גבוהה יותר, שמוצדקת בפרק~\ref{ch:budget}.
\end{keybox}

מהחלופות שנדחו נלקחו רכיבים: מחלופה ב׳, מתן ״מסלול ירוק״ מהיר לתמונות חתומות. מחלופה ג׳, אפשרות להסלים מקרים קשים לגוף בדיקה חיצוני בתשלום לפי שימוש.
